\chapter{Who This Book Is For}
\label{ch:who-for}

This book is written for a working data engineer or analytics engineer, two to
six years in, who runs Spark SQL on an Iceberg lakehouse most days and has a
nagging sense that the optimizer and the table format are black boxes. You may
recognize some of this:

\begin{itemize}
  \item You own pipelines that are fine until they are not --- a job that doubled
  in runtime after the data grew, a join that spills to disk, a query that scans
  the whole table when it should scan a day.
  \item You have opened the Spark UI, but you read it like a dashboard rather than
  a diagnostic.
  \item You treat \texttt{spark.sql.*} settings as spells found on a forum, and you
  could not say which mechanism a given one touches.
  \item You know the JVM exists and that garbage collection and memory matter,
  without a working model of either.
  \item You think of Iceberg as ``Parquet with time travel.''
\end{itemize}

If that describes you, the book is aimed at the gap between where you are and
where \Cref{tab:who-levels} says you can be.

\begin{table}[htbp]
\centering
\caption{Where the book starts you and where it takes you}
\label{tab:who-levels}
\small
\begin{tabularx}{\linewidth}{@{}l>{\raggedright\arraybackslash}X>{\raggedright\arraybackslash}X@{}}
\toprule
\textbf{Topic} & \textbf{Today} & \textbf{After the book} \\
\midrule
Query plans & ``There is an \texttt{EXPLAIN} somewhere.'' & Predicts the physical plan and its shuffles before running; reads \texttt{EXPLAIN FORMATTED} fluently. \\
Slowness & Guesses, adds memory, retries. & Localizes the cost to a stage and operator from the UI and metrics. \\
The shuffle & ``The slow part.'' & Knows the write and read paths, sizes partitions, removes unnecessary exchanges. \\
Joins & Lets Spark decide. & Predicts the strategy, forces it when statistics are wrong, fixes skew. \\
Iceberg & ``Parquet with snapshots.'' & Reads the metadata tree and predicts what a query will prune. \\
Configuration & Copies snippets. & Knows which knob touches which mechanism, and what it costs. \\
\bottomrule
\end{tabularx}
\end{table}

\section*{What you should already know}

The book assumes, and does not teach:

\begin{itemize}
  \item \textbf{SQL}: joins, aggregation, window functions, subqueries.
  \item \textbf{The relational model} and enough relational algebra to read a plan
  as a tree of operators.
  \item \textbf{Columnar versus row storage}, at the level of knowing what Parquet
  is for. \Cref{app:refenv} adds the background Part~III needs.
  \item \textbf{Kubernetes as an operator}: pods, resource limits, and why a pod is
  OOM-killed.
  \item Reading a stack trace, and basic command-line and Git use.
\end{itemize}

You do \emph{not} need to know Scala or Java. A few examples use Scala where SQL
cannot express the idea, and they are short and explained line by line.

\section*{What this book will not teach}

To keep the focus on query performance on a lakehouse, the book leaves out:

\begin{itemize}
  \item Scala or Java as languages;
  \item MLlib, GraphX, and machine-learning workloads;
  \item provisioning a cluster, or administering Kubernetes beyond what the
  executor sizing chapters require;
  \item the pandas API on Spark;
  \item Delta Lake and Apache Hudi, except where a comparison with Iceberg
  clarifies a design choice.
\end{itemize}

\section*{Who might want a different book}

If you are brand new to Spark and have never run a query on it, begin with an
introductory text and come back at \Cref{ch:shape}; the chapter is gentle, but it
does not teach what a DataFrame is. If you administer the cluster rather than
query it, Parts~I and~IV will serve you, and Part~III will explain what your
users' tables are doing to your object store. And if you already read query plans
for a living, the Iceberg chapters (Part~III) and the measured experiments are
where this book is most likely to tell you something new.
